\documentclass[12pt, a4paper]{article}
\usepackage{xeCJK} % Support for Chinese characters
\setCJKmainfont{Microsoft JhengHei}[
    AutoFakeBold=2.5,
    AutoFakeItalic=true
]
\usepackage{geometry}
\geometry{left=2.5cm, right=2.5cm, top=2.5cm, bottom=2.5cm}
\usepackage{booktabs}
\usepackage{float}

\begin{document}

\noindent
\textbf{個人報告} \hfill 系級：資工二 \quad 學號：113590017 \quad 姓名：黃楷程

\vspace{1.5em}
\begin{center}
    {\Large \textbf{個人心得報告}}
\end{center}
\vspace{1em}

本次 Lab 09 實驗中，我主要負責小組報告撰寫、簡報製作與技術規格說明的彙整。為了將 CPU 的架構與運作邏輯以最具條理的方式呈現，我仔細研讀了承諭編寫的 VHDL 核心程式碼，分析其硬體描述的實作細節，並將其轉化為具體的系統規格文件。

在報告製作的過程中，我系統化地剖析了 CPU 的資料流向。我詳細說明了 \texttt{LOAD Rs} 指令與 \texttt{MOVE Rs, Rt} 指令在硬體上的運作機制，並針對本次新增的算術指令（\texttt{ADD}、\texttt{AND}、\texttt{SLT}、\texttt{SUB A-B}、\texttt{DIV}、\texttt{SUB B-A}）進行深入的資料路徑分析。以除法（\texttt{DIV}）指令為例，資料經由自訂的 \texttt{udiv8} 函數進行 procedural 還原除法模擬，計算出的商即時輸出至 \texttt{Bus\_val}，並在時脈正緣時寫入 Rs。我將這些複雜的運算邏輯與多路選擇器選取路徑繪製成清晰的資料路徑圖，並在報告中對暫存器檔案的讀寫機制、匯流排的選擇邏輯進行了說明。

此外，我也詳細記錄了七段顯示器解碼模組（\texttt{hex\_to\_7seg}）與 FPGA 各輸出埠的對照關係（如 HEX0--HEX1 顯示 Bus 數值，HEX2--HEX5 顯示暫存器狀態）。這次的報告製作不僅培養了我技術文件的撰寫與架構視覺化的能力，更迫使我從宏觀的系統層面去理解處理器中控制單元與資料路徑之間的協同運作。這對於我理解計算機組織結構有著極大的助益，也為未來的專案合作奠定了良好的基礎。

\vspace{2em}
\noindent
\textbf{工作分配表}
\vspace{0.5em}

\begin{table}[H]
    \centering
    \begin{tabular}{llcl}
        \toprule
        \textbf{學號} & \textbf{姓名} & \textbf{比重} & \textbf{工作項目} \\ \midrule
        113590051 & 楊承諭 & 65\% & 小組專案製作與測試 \\
        113590017 & 黃楷程 & 35\% & 小組報告製作 \\ \bottomrule
    \end{tabular}
\end{table}

\end{document}
